% !TeX root = main.tex
% !TeX program = xelatex
% Discussion follows the evidence and claim scope in Chapters 3--4.

\section{讨论与结论}
\label{sec:discussion-conclusion}

SAGA 的分阶段授权保证依赖策略语义、凭证使用关系与并发状态更新的实际执行。
规则优先偏差可传播至 ACT 建立，OTK 转交反例在发放上限成立时违反增强的同人使用对应，ACT 分析则定位了检查与扣减分离造成的超额准入。

各阶段需要建立与目标性质一致的执行约束。
策略判定应按确定的优先语义选择规则，使最高优先匹配决定实际预算。
若以 OTK 领取预算进一步限制每个身份取得新 ACT 的次数，还需明确材料共享时的额度归属及实际使用身份的约束。
ACT 的每次准入应对应一个尚未被其他请求消耗的额度单位，检查与消耗须在并发执行中保持一致。
相应改进的实现与安全性仍需独立验证。

参考语义对照与受控执行将策略偏差追踪至 ACT 签发与存储，Tamarin 固定配置中的反例与对照刻画了发放和使用的身份关系，TLC 反例与原生方法实验分别观察到超额准入和额外后端调用。
各项结论限定于相应模型假设、资源配置及实验观察终点。
SAGA 的授权保证需要沿实际决策、凭证使用和请求准入的执行过程逐项核验。
